%========================================================
% D5 — Composite Phase--Depolarization Fiber Channel
%========================================================

\section{Composite Phase--Depolarization Fiber Channel}

This derivation document develops the analytical structure of the composite phase--depolarization fiber model introduced in Chapter~5. The objective is to derive explicitly the correlation tensor, purity, Bell-state fidelity, concurrence, and CHSH behavior associated with the simultaneous action of coherent phase evolution and local isotropic depolarization.

\medskip

%--------------------------------------------------------
% D5.1 — Composite Channel Definition
%--------------------------------------------------------

\subsection{Composite Channel Definition}

The effective composite channel is defined by

\begin{equation}
\rho_{\mathrm{out}}
=
\left(
\mathcal{D}_{p_A}
\otimes
\mathcal{D}_{p_B}
\right)
\left[
U_{\theta}
\rho_{\mathrm{in}}
U_{\theta}^{\dagger}
\right],
\label{eq:d5_composite_channel}
\end{equation}

where

\begin{equation}
U_{\theta}
=
E_{\theta}
\otimes
I,
\end{equation}

with

\begin{equation}
E_{\theta}
=
\begin{pmatrix}
1 & 0 \\
0 & e^{i\theta}
\end{pmatrix}.
\end{equation}

\medskip

The phase-evolved Bell state is therefore

\begin{equation}
|\psi(\theta)\rangle
=
\frac{1}{\sqrt{2}}
\left(
|01\rangle
+
e^{i\theta}|10\rangle
\right),
\label{eq:d5_phase_bell_state}
\end{equation}

with corresponding density operator

\begin{equation}
\rho(\theta)
=
|\psi(\theta)\rangle
\langle
\psi(\theta)
|.
\end{equation}

\medskip

As established in Appendix~C.2, isotropic local depolarization commutes with the reduced phase-evolution model:

\begin{equation}
\left(
\mathcal{D}_{p_A}
\otimes
\mathcal{D}_{p_B}
\right)
\left[
U_{\theta}
\rho
U_{\theta}^{\dagger}
\right]
=
U_{\theta}
\left[
\left(
\mathcal{D}_{p_A}
\otimes
\mathcal{D}_{p_B}
\right)
(\rho)
\right]
U_{\theta}^{\dagger}.
\label{eq:d5_channel_commutation}
\end{equation}

Therefore, coherent phase evolution and isotropic local depolarization may be treated independently while remaining equivalent under composition.

\medskip

Define the correlation survival factor

\begin{equation}
\eta
=
(1-p_A)(1-p_B).
\label{eq:d5_eta_definition}
\end{equation}

\medskip

%--------------------------------------------------------
% D5.2 — Correlation Tensor Derivation
%--------------------------------------------------------

\subsection{Correlation Tensor Derivation}

For the phase-evolved Bell state of Eq.~\eqref{eq:d5_phase_bell_state}, the correlation tensor is

\begin{equation}
T(\theta)
=
\begin{pmatrix}
\cos\theta & -\sin\theta & 0 \\
\sin\theta & \cos\theta & 0 \\
0 & 0 & -1
\end{pmatrix}.
\label{eq:d5_phase_tensor}
\end{equation}

\medskip

From the tensor-contraction law established in Chapter~5,

\begin{equation}
T_{\mathrm{out}}
=
\eta
T_{\mathrm{in}},
\end{equation}

one obtains

\begin{equation}
T_{\mathrm{out}}(\theta,p_A,p_B)
=
\eta
\begin{pmatrix}
\cos\theta & -\sin\theta & 0 \\
\sin\theta & \cos\theta & 0 \\
0 & 0 & -1
\end{pmatrix}.
\label{eq:d5_output_tensor}
\end{equation}

Therefore,

\begin{align}
T_{xx}^{\mathrm{out}}
&=
\eta\cos\theta,
\\
T_{xy}^{\mathrm{out}}
&=
-\eta\sin\theta,
\\
T_{yx}^{\mathrm{out}}
&=
\eta\sin\theta,
\\
T_{yy}^{\mathrm{out}}
&=
\eta\cos\theta,
\\
T_{zz}^{\mathrm{out}}
&=
-\eta.
\end{align}

All remaining tensor components vanish.

\medskip

Thus, coherent phase evolution rotates the transverse \(xy\)-correlation structure, while local depolarization contracts the tensor magnitude uniformly.

\medskip

%--------------------------------------------------------
% D5.3 — Purity Derivation
%--------------------------------------------------------

\subsection{Purity Derivation}

For Bell-state inputs, the reduced local Bloch vectors vanish. The output state may therefore be written as

\begin{equation}
\rho_{\mathrm{out}}
=
\frac{1}{4}
\left[
I\otimes I
+
\sum_{i,j=1}^{3}
T_{ij}^{\mathrm{out}}
\sigma_i\otimes\sigma_j
\right].
\label{eq:d5_pauli_density_expansion}
\end{equation}

The global purity is

\begin{equation}
\gamma_{AB}
=
\mathrm{Tr}
\left(
\rho_{\mathrm{out}}^{2}
\right).
\end{equation}

Substituting Eq.~\eqref{eq:d5_pauli_density_expansion} gives

\begin{align}
\gamma_{AB}
&=
\frac{1}{16}
\mathrm{Tr}
\left[
\left(
I\otimes I
+
\sum_{i,j}T_{ij}\sigma_i\otimes\sigma_j
\right)^2
\right].
\end{align}

Using the orthogonality relation

\begin{equation}
\mathrm{Tr}
\left[
(\sigma_i\otimes\sigma_j)
(\sigma_k\otimes\sigma_l)
\right]
=
4
\delta_{ik}
\delta_{jl},
\label{eq:d5_pauli_orthogonality}
\end{equation}

one obtains

\begin{equation}
\gamma_{AB}
=
\frac14
\left[
1
+
\sum_{i,j=1}^{3}
(T_{ij}^{\mathrm{out}})^2
\right].
\label{eq:d5_purity_tensor_formula}
\end{equation}

\medskip

From Eq.~\eqref{eq:d5_output_tensor},

\begin{align}
\sum_{i,j=1}^{3}
(T_{ij}^{\mathrm{out}})^2
&=
\eta^2
\left[
\cos^2\theta
+
\sin^2\theta
+
\sin^2\theta
+
\cos^2\theta
+
1
\right]
\nonumber
\\
&=
3\eta^2.
\end{align}

Therefore,

\begin{equation}
\boxed{
\gamma_{AB}
=
\frac{1+3\eta^2}{4}
}
\label{eq:d5_final_purity}
\end{equation}

with

\[
\eta=(1-p_A)(1-p_B).
\]

\medskip

Since the reduced states remain maximally mixed,

\begin{equation}
\rho_A
=
\rho_B
=
\frac{I}{2},
\end{equation}

their purities are

\begin{equation}
\boxed{
\gamma_A
=
\gamma_B
=
\frac12
}.
\end{equation}

\medskip

%--------------------------------------------------------
% D5.4 — Bell-State Fidelity Derivation
%--------------------------------------------------------

\subsection{Bell-State Fidelity Derivation}

The Bell-state fidelity relative to the reference state

\[
|\Psi^+\rangle
=
\frac1{\sqrt2}
(|01\rangle+|10\rangle)
\]

is

\begin{equation}
F_{\Psi^+}
=
\langle\Psi^+|
\rho_{\mathrm{out}}
|\Psi^+\rangle.
\end{equation}

\medskip

For the pure phase-evolved state,

\begin{align}
\langle\Psi^+|\psi(\theta)\rangle
&=
\frac12
\left(
1+e^{i\theta}
\right),
\end{align}

so that

\begin{align}
F_{\Psi^+}^{\mathrm{phase}}
&=
\left|
\frac12
(1+e^{i\theta})
\right|^2
\nonumber
\\
&=
\frac14
(2+2\cos\theta)
\nonumber
\\
&=
\frac{1+\cos\theta}{2}.
\label{eq:d5_phase_fidelity}
\end{align}

\medskip

Under local depolarization, the bipartite coherences are contracted by the factor \( \eta \). Therefore,

\begin{equation}
\boxed{
F_{\Psi^+}
=
\frac{1+\eta\cos\theta}{2}
}
\label{eq:d5_final_fidelity}
\end{equation}

with

\[
\eta=(1-p_A)(1-p_B).
\]

\medskip

This expression separates coherent phase mismatch from incoherent correlation degradation.

\medskip

%--------------------------------------------------------
% D5.5 — Concurrence Derivation
%--------------------------------------------------------

\subsection{Concurrence Derivation}

The composite channel produces a locally unitary-equivalent Bell-diagonal family. Since local unitary transformations preserve concurrence, the phase parameter \( \theta \) does not affect the entanglement measure.

\medskip

The concurrence therefore depends only on the correlation survival factor \( \eta \).

\medskip

For Bell-diagonal states with isotropically contracted correlations, one obtains

\begin{equation}
\boxed{
C
=
\max
\left(
0,
\frac{3\eta-1}{2}
\right)
}
\label{eq:d5_final_concurrence}
\end{equation}

with

\[
\eta=(1-p_A)(1-p_B).
\]

\medskip

Entanglement disappears when

\begin{equation}
\frac{3\eta-1}{2}=0,
\end{equation}

which gives

\begin{equation}
\boxed{
\eta=\frac13
}.
\label{eq:d5_entanglement_threshold}
\end{equation}

\medskip

For symmetric local depolarization,

\[
p_A=p_B=p,
\]

one obtains

\begin{equation}
(1-p)^2
=
\frac13,
\end{equation}

and therefore

\begin{equation}
\boxed{
p
=
1-\frac1{\sqrt3}
\approx
0.423
}.
\end{equation}

\medskip

%--------------------------------------------------------
% D5.6 — CHSH Scaling Derivation
%--------------------------------------------------------

\subsection{CHSH Scaling Derivation}

The maximal CHSH parameter is determined through the Horodecki criterion:

\begin{equation}
S_{\max}
=
2\sqrt{\lambda_1+\lambda_2},
\label{eq:d5_horodecki}
\end{equation}

where \( \lambda_1,\lambda_2 \) are the two largest eigenvalues of

\begin{equation}
T^TT.
\end{equation}

\medskip

Since

\begin{equation}
T_{\mathrm{out}}
=
\eta T_{\mathrm{in}},
\end{equation}

one obtains

\begin{align}
T_{\mathrm{out}}^T
T_{\mathrm{out}}
&=
(\eta T_{\mathrm{in}})^T
(\eta T_{\mathrm{in}})
\nonumber
\\
&=
\eta^2
T_{\mathrm{in}}^T
T_{\mathrm{in}}.
\end{align}

Therefore, all eigenvalues scale according to

\begin{equation}
\lambda_k^{\mathrm{out}}
=
\eta^2
\lambda_k^{\mathrm{in}}.
\end{equation}

Substituting into Eq.~\eqref{eq:d5_horodecki} gives

\begin{align}
S_{\max}^{\mathrm{out}}
&=
2\sqrt{
\eta^2
(\lambda_1^{\mathrm{in}}+\lambda_2^{\mathrm{in}})
}
\nonumber
\\
&=
\eta
S_{\max}^{\mathrm{in}}.
\end{align}

Hence,

\begin{equation}
\boxed{
S_{\max}^{\mathrm{out}}
=
\eta
S_{\max}^{\mathrm{in}}
}
\label{eq:d5_final_chsh_scaling}
\end{equation}

with

\[
\eta=(1-p_A)(1-p_B).
\]

\medskip

For Bell states,

\[
S_{\max}^{\mathrm{in}}
=
2\sqrt2,
\]

therefore

\begin{equation}
S_{\max}^{\mathrm{out}}
=
2\sqrt2\,\eta.
\end{equation}

Bell-inequality violation requires

\begin{equation}
S_{\max}^{\mathrm{out}}>2,
\end{equation}

which gives

\begin{equation}
\boxed{
\eta>\frac1{\sqrt2}
}.
\label{eq:d5_chsh_threshold}
\end{equation}

\medskip

For symmetric local depolarization,

\begin{equation}
(1-p)^2
>
\frac1{\sqrt2},
\end{equation}

thus

\begin{equation}
\boxed{
p
<
1-2^{-1/4}
\approx
0.159
}.
\end{equation}

\medskip

Therefore, Bell nonlocality disappears before entanglement vanishes completely, showing that CHSH violation is more fragile than concurrence under local isotropic depolarization.